(1) 导轨几何给 \(\dot y=2ax\dot x\)，因此 \(L=\frac12m(1+4a^2x^2)\dot x^2-mga x^2\)。按变质量度量求导可得 \((1+4a^2x^2)\ddot x+4a^2x\dot x^2+2ga x=0\)；直接把水平分量当作直线简谐运动会遗漏曲线度量项。
(2) 初态能量 \(E=mgaA^2\)，故 \(\dot x^2=2ga(A^2-x^2)/(1+4a^2x^2)\)。沿一象限积分并乘四，得到 \(T=4/\sqrt{2ga}\int_0^{\pi/2}\sqrt{1+4a^2A^2\sin^2u}\,du\)，该式对有限 \(A\) 也成立，不需把弧长误认为水平位移。
(3) 展开平方根 \(\sqrt{1+4a^2A^2\sin^2u}=1+2a^2A^2\sin^2u+O((aA)^4)\)；利用 \(\int_0^{\pi/2}\sin^2u\,du=\pi/4\) 得 \(T=(2\pi/\sqrt{2ga})[1+a^2A^2+O((aA)^4)]\)。大幅度时每单位水平位移的实际弧长增加，因此周期的领先修正为正。
(4) 取向上指向曲率中心的单位法向 \(\mathbf n=(-2ax,1)/\sqrt{1+4a^2x^2}\)。珠子加速度在该向的分量为 \(2a\dot x^2/\sqrt{1+4a^2x^2}\)，所以 \(N=m(g+2a\dot x^2)/\sqrt{1+4a^2x^2}>0\)。在谷底 \(x=0\)，\(\dot x^2=2gaA^2\)，故 \(N=mg(1+4a^2A^2)\)；在转折点 \(x=\pm A\)，\(N=mg/\sqrt{1+4a^2A^2}\)。
